\section*{Bab \ref{ch:b3:identical-particles} --- Partikel Identik}

\begin{solution}{exo:b3:identical-particles:1}
(a) Menukar dua kali sama dengan tidak berbuat apa pun: \omterm{def:b3:quantum-formalism:observable}{nilai eigennya}
$\pm1$. (b) \omterm{def:b3:identical-particles:exchange}{Partikel identik} masuk ke $\hat H$ secara setangkup,
sehingga $\hat H\hat
P_{12} = \hat P_{12}\hat H$: golongan kesetangkupannya tak pernah
berubah. (c)
$\ket a\otimes\ket b$ dan $\ket b\otimes\ket a$ akan menjadi keadaan yang
\emph{berbeda} padahal memerikan keadaan fisis yang sama --- yakni
pembedaan tak teramati yang tak boleh dibawa formalismenya. (d)
$^1$H (2 fermion): boson; deuterium (3): fermion; $^4$He: boson;
$^3$He: fermion; H$_2$ (4): boson; foton: boson.
\end{solution}

\begin{solution}{exo:b3:identical-particles:2}
(a) $(\ket a\otimes\ket b - \ket b\otimes\ket a)/\sqrt2$. (b)
Nol secara identik. (c) Jumlahan berselang-seling atas keenam
permutasinya --- yakni determinan $3\times3$ orbital terhadap
partikelnya; menukar dua barisnya membalik tanda determinannya. (d)
Determinan dengan dua kolom yang sama akan lenyap: dua fermion dalam
satu orbital adalah kolom yang terulang.
\end{solution}

\begin{solution}{exo:b3:identical-particles:3}
(a) $3^2 = 9$. (b) $3$ pasangan ganda $+$ $3$ pasangan campuran $= 6$.
(c) $3$ pasangan.
(d) Enam keadaan satu partikel (tiga orbital $\times$ dua spin):
$\binom62 = 15$.
\end{solution}

\begin{solution}{exo:b3:identical-particles:4}
(a) C: $1s^22s^22p^2$; N: $2p^3$; O: $2p^4$; Ne: $2p^6$. (b)
Ketiga elektron $2p$ nitrogen menempati ketiga orbitalnya
sendiri-sendiri, dengan spin sejajar. (c) Spin yang sejajar memaksa
keadaan ruang yang antisetangkup, yang lubang penukarannya menghemat
tolakan Coulomb --- penyejajarannya dibeli dengan penghematan
elektrostatik. (d) Helium (beserta kerabat mulianya); fluorin, yang
kurang satu elektron dari penutupan neon, merenggutnya dari hampir apa
saja.
\end{solution}

\begin{solution}{exo:b3:identical-particles:5}
(a) Boson: $2E_1$; fermion tanpa spin: $E_1 + 4E_1 = 5E_1$;
elektron: $2E_1$ (spin yang berlawanan berbagi $n = 1$). (b) Boson dan
elektron: $E_1 + 4E_1 = 5E_1$; fermion tanpa spin: $E_1 + 9E_1 =
10E_1$. (c) Memampatkan fermion memaksa pendudukan aras yang makin
tinggi: ongkos energi tiap perasannya adalah tekanan ke luar yang tak
dapat dihilangkan pendinginan mana pun --- yakni ketaktermampatan materi
dalam bentuk embrio. (d)
$\sum_{n=1}^N n^2 \approx N^3/3$: total lautan Ferminya tumbuh sebagai
$N^3$.
\end{solution}

\begin{solution}{exo:b3:identical-particles:6}
(a) Kedua sukunya saling meniadakan pada $\vect r_1 = \vect r_2$. (b)
Setelah dijabarkan, suku silangnya menyumbang $\mp2|\bra a\hat{\vect r}\ket b|^2$:
keadaan antisetangkupnya mempunyai rata-rata kuadrat pemisahan yang
\emph{lebih besar}. (c) Lebih berjauhan berarti energi tolakannya lebih
kecil: $E_- <
E_+$. (d) Triplet helium (ruang antisetangkup) terletak
di bawah singlet konfigurasi yang sama --- ortohelium di bawah
parahelium, sebesar \qty{0.8}{eV} tolakan yang dihemat.
\end{solution}

\begin{solution}{exo:b3:identical-particles:7}
(a) $1s^2$ adalah keadaan ruang setangkup: spinnya harus membentuk
singlet yang antisetangkup. (b) Pembelahannya $= 2J$: $J \approx
\qty{0.4}{eV}$. Energi dipol--dipol magnetik yang sejati pada
$a_0$: $\mu_0\mu_{\text{B}}^2/4\pi a_0^3 \approx \qty{4e-4}{eV}$
--- seribu kali lebih kecil: ``gaya spin'' itu adalah elektrostatika
yang mengenakan warna spin. (c) Meluruh ke $1\,^1S$ menuntut pembalikan
spin (singlet--triplet) \emph{dan} transisi $s \to s$ ($\Delta l =
0$): jadi terlarang ganda. (d) $8000$ s berbanding \qty{1.6}{ns}: dua
belas orde besar, dari dua kaidah pemilihan --- kemetastabilan adalah
kaidah yang bertumpuk, bukan gaya yang melemah.
\end{solution}

\begin{solution}{exo:b3:identical-particles:8}
(a) Tanda penukaran totalnya $=$ (bagian spinnya) $\times (-1)^J$
haruslah $-1$: triplet (spin setangkup) mengambil $J$ ganjil, dan
singletnya $J$ genap.
(b) $3:1$ dari kelipatgandaannya. (c) Bagian orto sebesar 3/4 mengendur
menjadi para ($J = 0$), sambil melepaskan $2B \approx \qty{15}{meV}$
tiap molekul --- \emph{lebih besar} daripada kalor laten zat cairnya:
sebuah tangki tanpa mangkin perlahan mendidihkan dirinya sampai kering.
Pabriknya mengubah orto $\to$ para di atas mangkin selama pencairannya.
(d) Peroketan hidrogen cair --- tiap landasan peluncuran mewarisi
persoalan statistik spin inti.
\end{solution}

\begin{solution}{exo:b3:identical-particles:9}
(a) Lajunya berbanding kuadrat amplitudonya: $n + 1$ ke dalam ragam yang
terisi, dan $n$ untuk arah sebaliknya. (b) $10^{10} + 1$ berbanding $1$:
ragam yang ramai menang sepuluh miliar kali. (c) Satu foton spontan
dalam ragam yang disukai memiringkan peluangnya; tiap foton terangsang
menaikkan $n$ lalu mempercuram kemiringannya: yakni longsoran foton yang
identik --- penguatan yang menyempit ke dalam satu ragam tunggal. (d)
Dalam $\hat a^\dagger\ket n =
\sqrt{n+1}\,\ket{n+1}$: ragam bosoniknya adalah osilatornya, dan
penguatannya adalah koefisien tangganya.
\end{solution}

\begin{solution}{exo:b3:identical-particles:10}
(a) Tiap gasnya mengembang ke volume dua kali lipat: $\Delta S =
Nk_{\text{B}}\ln2$ tiap gas. (b) Menutup lalu membuka kerannya lagi
akan memproduksi entropi tanpa perubahan keadaan apa pun --- yakni
kecurangan pembukuan abadi. (c) Dengan molekul yang bernama, gas
gabungannya mempunyai $\binom{2N}{N}$ penetapan ``di sisi mana''
tambahan yang bernilai $2Nk_{\text{B}}\ln2$; membagi semua cacah
keadaannya dengan $N!$ --- yakni pencacahan jujur bagi \omterm{def:b3:identical-particles:exchange}{partikel identik}
--- menyingkirkan persis itu.
(d) $N!$ yang diperlukan fisika klasik namun tak dapat dibenarkannya
adalah postulat penyetangkupan yang dilihat dari kejauhan.
\end{solution}

\begin{solution}{exo:b3:identical-particles:11}
(a) $E_{\text{F}} = \hbar^2(3\pi^2n_{\text{e}})^{2/3}/2m_{\text{e}}
\approx \qty{0.16}{MeV}$ --- sepertiga \omterm{def:b3:relativistic-dynamics:momentum-energy}{energi diam} elektronnya:
bintangnya duduk di tepi rezim relativistik. (b)
$k_{\text{B}}T \approx \qty{0.9}{keV} \ll E_{\text{F}}$: pada skala
Ferminya bintang yang berpijar putih itu sangat berdegenerasi --- jadi
``dingin''.
(c) Elektronnya menumpuk secara Pauli sampai
$\qty{0.16}{MeV}$; memampatkan bintangnya menaikkan tiap aras yang
terisi, dan gradien energi itu adalah tekanan yang tidak bergantung
suhu --- asas larangannya, bukan kalornya, yang memikul bebannya. (d)
Degenerasi relativistiknya melunakkan hukum tekanannya sampai
gravitasinya menang pada massa yang berhingga: yakni batas
Chandrasekhar, yang disimpan untuk \cref{ch:b3:astrophysics}.
\end{solution}

\begin{solution}{exo:b3:identical-particles:12}
(a) $\hat{\vect S}_1\cdot\hat{\vect S}_2 = \tfrac12(\hat S^2 -
\tfrac32\hbar^2)$: $-\tfrac34\hbar^2$ pada singletnya,
$+\tfrac14\hbar^2$ pada tripletnya --- \omterm{def:b3:hamiltonian-mechanics:hamiltonian}{Hamiltonian} yang dinyatakan itu
menghasilkan kembali pembelahan $2J$-nya. (b) Sebuah kisi berisi
tetangga yang menyearah: kemagnetan spontan --- yakni feromagnetisme.
(c) $J/k_{\text{B}}
\approx \qty{1200}{K}$: nilai \qty{1043}{K} milik besi tepat pada skala
ini. (d) Kemagnetan dipol--dipol berorde $10^{-4}$ eV: ia akan menata
diri pada milikelvin. Feromagnetisme adalah tolakan Coulomb yang memilih
susunan spinnya lewat geometri penukarannya.
\end{solution}

\begin{solution}{pb:b3:identical-particles:1}
\textbf{1.} Orbital $(n, l, m)$ dengan $n^2$ tiap $n$; spinnya
melipatduakan; Pauli membatasi satu elektron tiap label lengkapnya:
$2$, $8$, $18$.
\textbf{2.} Elektron dalamnya menyaring intinya secara berbeda bagi
$l$ yang berbeda: orbital $s$ yang menembus tenggelam di bawah orbital
$p$ dan $d$ yang kurang menembus --- urutan energinya mengikuti
$(n, l)$, bukan $n$ saja.
\textbf{3.} Urutannya $1s\,2s\,2p\,3s\,3p\,4s\,3d$; Na: [Ne]$3s^1$;
Cl: [Ne]$3s^23p^5$; K: [Ar]$4s^1$; Fe: [Ar]$4s^23d^6$.
\textbf{4.} Barisnya menutup pada tiap konfigurasi mulianya: $2$ ($1s$),
$8$ ($2s2p$), $8$ lagi ($3s3p$ --- karena $3d$ terletak di atas $4s$
dan menunggu), lalu $18$ ($4s3d4p$).
\textbf{5.} Sebuah kulit tertutup: celahnya besar menuju orbital
berikutnya, tak ada elektron murah untuk diberikan, dan tak ada ruang
untuk menerima.
\textbf{6.} $Z_{\text{eff}}$ adalah muatan intinya dikurangi
penyaringan elektron dalamnya (dan sebagian, elektron sekulitnya):
elektron valensi natrium melihat $11 - 10 \approx 1$; kedua elektron
helium berbagi $Z = 2$ yang telanjang, dengan masing-masing menyaring
temannya hanya sepertiga.
\textbf{7.} Li $\to$ berilium naik, lalu turun pada B; N $\to$ O turun:
itulah kedua lekukannya.
\textbf{8.} Boron membuka $2p$, yang lebih tinggi dan kurang menembus
daripada $2s$: elektron barunya memang lebih mudah dibuang.
\textbf{9.} Oksigen harus, untuk pertama kalinya, \emph{memasangkan}
dua elektron dalam satu orbital $p$: tolakannya bertambah, dan
pemantapan penukarannya hilang --- elektron $2p$ keempatnya mendapat
potongan harga.
\textbf{10.} Tiap elektronnya menyaring temannya secara tak sempurna:
pencocokan $24.6 = 13.6\,Z_{\text{eff}}^2$ memberikan $Z_{\text{eff}} =
1.34$ --- dua pertiga muatan yang tersaring.
\textbf{11.} Perisai terasnya yang nyaris lengkap ($Z_{\text{eff}}
\approx 1.3$) dan elektron terluarnya yang ber-$n = 2$: keduanya
memangkas ikatannya.
\textbf{12.} Elektron yang ditambahkan masuk ke kulit yang \emph{sama}
lalu menyaring satu sama lain dengan buruk sementara $Z$ memanjat:
$Z_{\text{eff}}$ tumbuh dan seluruh kulitnya menciut --- neon lebih
kecil daripada litium.
\textbf{13.} Satu elektron $s$ yang murah di atas teras mulia:
pengionannya $\sim\qty{5}{eV}$, ion $+1$ seketika, dan penyerahan
elektron yang ganas --- ciri keluarganya hanyalah satu konfigurasi.
\textbf{14.} Satu lubang dalam penutupan $p^6$: afinitas elektronnya
besar, ionnya $-1$, dan lubangnya berpasangan menjadi molekul X$_2$.
\textbf{15.} $2s$ dan $2p$ terletak berdekatan: menaikkan
$2s^22p^2 \to$ empat hibrida setengah terisi berongkos sedikit namun
membeli dua ikatan tambahan --- kevalensiempatan karbon adalah
kenyaris-degenerasian yang dimanfaatkan.
\textbf{16.} Dengan empat elektron, gabungan pengikat \emph{maupun}
pengantiikatnya sama-sama terisi: perolehannya saling meniadakan,
sehingga tak ada ikatan kovalen --- helium tetap penyendiri.
\textbf{17.} Elektron $d$-nya mengubur diri di dalam kulit $4s$-nya:
kimia luarnya, yang ditetapkan kulit $s$-nya, sedikit sekali berubah
sepanjang sepuluh unsur --- yakni dataran kimia.
\textbf{18.} Skandium, [Ar]$4s^23d^1$: elektron $d$ tak berpasangan
yang pertama.
\textbf{19.} Sebuah momen magnetik dari tiga spin yang menyearah
($\sim3\mu_{\text{B}}$); penyejajarannya dibayar penukarannya ---
yakni penghematan Coulomb, bukan gaya magnetik.
\textbf{20.} Elektron yang sedikit dapat dibedakan dapat berkerumun
menuju $1s$ (tanpa peniadaan eksak keadaan tersetangkupkannya):
kulitnya akan bocor, penutupannya mengabur, dan hukum berkalanya larut.
\textbf{21.} $v/c \approx Z\alpha = 0.58$: pertambahan massa
relativistiknya menciutkan dan memantapkan orbital $6s$-nya; serapan
$5d \to 6s$ emas meluncur ke arah biru (sehingga menyisakan pantulan
kuning), dan pasangan $6s^2$ raksa yang mengetat berikatan begitu lemah
sehingga logamnya cair.
\textbf{22.} Keduanya berjalan nyaris sejajar: urutannya bergantung
pendudukan, $4s$ terionkan lebih dahulu, dan peminjaman $3d^5$/$3d^{10}$
oleh Cr dan Cu --- subkulit setengah terisi dan terisi penuh membeli
kemantapan penukaran dengan murah.
\textbf{23.} Kimia adalah urusan elektronnya, dan elektronnya hanya
merasakan \emph{muatan} intinya: massanya masuk sekadar lewat
pergeseran getaran (isotop) yang mungil.
\textbf{24.} Struktur kulitnya (yang memihak logika gas mulia)
berhadapan dengan pemantapan relativistik dan tegangan Coulomb yang
sesungguhnya pada $Z\alpha \to 1$: kimia superberat adalah tarik-menarik
keduanya.
\textbf{25.} Satu katalog orbital, satu pelipatduaan, satu tanda minus:
baris $2, 8, 8, 18$, lekukan pada boron dan oksigen, keluarga dari
alkali sampai mulia, dan tangga pengionan dari \qty{5.1}{eV} sampai
\qty{24.6}{eV} --- hukum berkala yang diturunkan, bukan dihafalkan.
\end{solution}
